\section{Логические правила и фразы}
\label{sec:work}

Ответ строит агент стандартного ответа: правила берутся по приоритету (сначала~1, затем~2 и~3). Сработавшее правило указывает класс фраз. Фраза выбирается случайно; переменные шаблона заменяются содержимым ссылок из БЗ.

Для каждого из шести предметных намерений заданы три правила. Более конкретная посылка имеет меньший номер приоритета, поэтому при полном фрагменте БЗ срабатывает подробный ответ, при неполном~\mdash  короткий или ответ <<неизвестно>>.

\begin{table}[H]
	\centering
	\small
	\begin{tabular}{|p{2.4cm}|p{3.4cm}|p{3.4cm}|p{3.0cm}|}
		\hline
		Намерение & Приоритет 1 & Приоритет 2 & Приоритет 3 \\
		\hline
		Исполнитель & исполнитель~\mdash  группа, известен участник & исполнитель указан & исполнитель не указан \\
		\hline
		Автор & автор принадлежит классу исполнителей & автор указан & автор не указан \\
		\hline
		Альбом & известны альбом и лейбл & известен только альбом & альбом не указан \\
		\hline
		Жанр & жанр и исполнитель & жанр указан & жанр не указан \\
		\hline
		Инструмент & инструмент и автор & инструмент указан & инструмент не указан \\
		\hline
		Год & альбом и лейбл (год выпуска) & год премьеры при известном авторе & год не указан \\
		\hline
	\end{tabular}
	\caption{Правила ответа и зависимость от состояния БЗ}
	\label{tab:rules}
\end{table}

Посылка каждого правила содержит \texttt{\_message <- concept\_message\_about\_\ldots} и дугу \texttt{\_message -> rrel\_entity: \_work}. Дальше проверяются отношения экземпляра: \texttt{nrel\_performer\_of}, \texttt{nrel\_author}, \texttt{nrel\_included\_in\_album}, \texttt{nrel\_work\_genre}, \texttt{nrel\_performed\_on\_instrument}, \texttt{nrel\_premiere\_year}, \texttt{nrel\_release\_year} и связанные факты (участник группы, лейбл).

Фрагмент правила <<автор сам является исполнителем>>:

\begin{lstlisting}[caption={Правило ответа, если автор входит в класс исполнителей},label={lst:rule-author},basicstyle=\footnotesize\ttfamily]
@if_author_performer = [*
	_message
		_<- concept_message;
		_<- concept_message_about_author;
		_-> rrel_entity:: _work;;
	_work _=> nrel_author:: _author;;
	_author _<- concept_performer;;
*];;

@then_author_performer = [*
	_message _=> nrel_reply:: _reply_message;;
	_reply_message
		_<- concept_message_about_author;
		_<- concept_atomic_message;;
*];;

lr_message_about_author_performer
<- concept_answer_on_standard_message_rule;
<- concept_answer_on_standard_message_rule_priority_1;
=> nrel_answer_pattern:
	{ rrel_1: concept_phrase_about_author_performer };
-> rrel_main_key_sc_element: @implication_author_performer;;
\end{lstlisting}

Заданы восемнадцать классов фраз, по два русских варианта в каждом (36 формулировок, без повторов). Все фразы шаблонные: подставляются имя произведения, исполнителя, участника группы, автора, альбома, лейбла, жанра, инструмента и года. Это больше требуемых трёх шаблонных фраз на диалог.

Примеры шаблонных фраз (подстановки~\mdash  имена из БЗ):
\begin{itemize}
	\item <<произведение исполняет группа~\ldots Один из её участников~\mdash  \ldots>>;
	\item <<\ldots написал \ldots>>;
	\item <<\ldots входит в альбом~\ldots (год), выпущенный лейблом~\ldots>>;
	\item <<\ldots относится к жанру~\ldots; исполняет~\ldots>>;
	\item <<\ldots исполняется на инструменте~\ldots>>;
	\item <<премьера~\ldots состоялась в~\ldots году>>.
\end{itemize}

Шаблон фразы повторяет посылку соответствующего правила и снимает \texttt{nrel\_main\_idtf} с найденных узлов. Если конструкция шаблона в памяти не находится, текст не генерируется и берётся следующая фраза класса.

\section{Тестирование}
\label{sec:test}

Одно и то же намерение даёт разные ответы, потому что в БЗ у произведений разный набор фактов: у Bohemian Rhapsody исполнитель~\mdash  группа Queen с участником Фредди Меркьюри; у Billie Jean исполнитель~\mdash  Майкл Джексон (не группа); у оперы <<Кармен>> исполнитель не указан. Аналогично для автора, альбома, жанра, инструмента и года.

\begin{table}[H]
	\centering
	\scriptsize
	\begin{tabular}{|p{4.2cm}|p{2.4cm}|p{6.0cm}|}
		\hline
		Вопрос & Правило & Ответ в интерфейсе NIKA \\
		\hline
		Кто исполняет Bohemian Rhapsody? & группа и участник & Bohemian Rhapsody исполняет группа Queen. Один из её участников~\mdash  Фредди Меркьюри. \\
		\hline
		Кто исполняет Billie Jean? & исполнитель & Исполнитель произведения Billie Jean~\mdash  Майкл Джексон. \\
		\hline
		Кто исполняет Кармен? & неизвестен & Я пока не знаю, кто исполняет Кармен. \\
		\hline
		Кто автор Billie Jean? & автор-исполнитель & Автор произведения Billie Jean~\mdash  исполнитель Майкл Джексон. \\
		\hline
		Кто автор Yesterday? & автор & Yesterday написал Пол Маккартни. \\
		\hline
		В какой альбом входит Come Together? & альбом и лейбл & Come Together входит в альбом Abbey Road (1969), выпущенный лейблом Apple Records. \\
		\hline
		Какой жанр у Bohemian Rhapsody? & жанр и исполнитель & Bohemian Rhapsody относится к жанру рок-музыка; исполняет Queen. \\
		\hline
		Какой жанр у Кармен? & жанр & Жанр произведения Кармен~\mdash  классическая музыка. \\
		\hline
		На каком инструменте исполняется Лунная соната? & инструмент & Лунная соната исполняется на инструменте фортепиано. \\
		\hline
		На каком инструменте исполняется Кармен? & неизвестен & Я не знаю, на каком инструменте исполняется Кармен. \\
		\hline
		В каком году премьера Лебединое озеро? & премьера & Премьера произведения Лебединое озеро состоялась в 1877 году. \\
		\hline
		В каком году вышла Come Together? & выпуск альбома & Come Together вышла в 1969 году в составе альбома Abbey Road. \\
		\hline
	\end{tabular}
	\caption{Разные ответы на одно намерение в зависимости от БЗ}
	\label{tab:tests}
\end{table}

Сценарий из восьми пар реплик, включая шаблонные ответы лабораторной работы №3, зафиксирован в sc-web (рисунок~\ref{fig:dialog}).

\begin{figure}[H]
	\centering
	\includegraphics[width=\textwidth]{dialog-scg}
	\caption{Диалог пользователя с Никой в базе знаний}
	\label{fig:dialog}
\end{figure}
